% ===== 原文 pp.57--60 =====

% [原文 p.57]

\subsubsection{一般的非阿贝尔四维规范理论}

现在考虑 3.3.5 节所考虑类型的一般 4d $\cG$ 规范理论。我们加入一个离散 theta 角，它由一个同态指定
\begin{equation}
    \gamma:\ \cZ \to \widehat{\cZ} \,.
    \tag{3.195}
\end{equation}
相应的理论被称为 $\cG_\gamma$ 规范理论。下面我们将讨论规范理论拉氏量中相应项的形式。

\para{'t Hooft--Wilson 线算符。}考虑一个诱导出 $\cG$ 的磁单极子位形的线算符，它具有如下性质：
\begin{equation}
    \int_{S^2} w_2 = z \in \cZ \,,
    \tag{3.196}
\end{equation}
其中 $S^2$ 环绕该线。这样的线算符必须按规范代数 $\fg$ 的不可约表示 $R$ 变换，且在 $\cZ$ 下的荷为
\begin{equation}
    q(R) = \gamma(z) \in \widehat{\cZ}
    \tag{3.197}
\end{equation}
注意，若 $\gamma(z) = 0$，则 $R$ 是规范群 $\cG$ 的表示；但若 $\gamma(z) \neq 0$，则 $R$ 不是 $\cG$ 的表示。

\para{1-形式对称性。}只计入磁屏蔽与场强导致的屏蔽，我们把未被屏蔽的线算符构成的群约化为
\begin{equation}
    \mathsf{D}_\gamma \subset \cZ \times \widehat{Z}(G) \,.
    \tag{3.198}
\end{equation}
它由元素
\begin{equation}
    (z, \tilde{\gamma}(z)) \in \cZ \times \widehat{Z}(G) \,,
    \tag{3.199}
\end{equation}
组成，其中 $z \in \cZ$，而 $\tilde{\gamma}(z) \in \widehat{Z}(G)$ 是 $\gamma(z) \in \widehat{\cZ}$ 在满同态
\begin{equation}
    \hat{i}:\ \widehat{Z}(G) \to \widehat{\cZ} \,,
    \tag{3.200}
\end{equation}
下的任意提升；该满同态是单同态
\begin{equation}
    i:\ \cZ \to Z(G) \,,
    \tag{3.201}
\end{equation}
的 Pontryagin 对偶，它对应于 $\cZ$ 是 $Z(G)$ 的子群这一事实。
代入 $z = 0$，我们看到纯 Wilson 线构成一个子群
\begin{equation}
    \widehat{Z}(\cG) \subseteq \mathsf{D}_\gamma \,.
    \tag{3.202}
\end{equation}

% [Tutor Note: p.57 式 (3.199)--(3.202) 的逻辑：\hat{i}: \hat{Z}(G) → \hat{\cZ} 是 i: Z → Z(G) 的对偶满同态，其核由在 i(Z) 上取 1 的特征标组成，即 Z(G)/Z 的特征标群 = \hat{Z}(\cG)（因为 Z(G)/Z 正是规范群 G 的中心）。把 z = 0 代入 (3.199)，\tilde{\gamma}(0) 跑遍这个核，于是得到 (3.202)。]

% [原文 p.58]

这来自以下事实：$\widehat{Z}(\cG) = \widehat{Z}(G)/\widehat{\cZ}$，因此平凡元素 $0 \in \widehat{\cZ}$ 在 $\widehat{Z}(G)$ 中的可能提升确实由 $\widehat{Z}(\cG)$ 的元素给出。

物质内容完全屏蔽了这个子群的一部分
\begin{equation}
    Q_{\mathrm{matter}} \subseteq \widehat{Z}(\cG) \subseteq \mathsf{D}_\gamma \,,
    \tag{3.203}
\end{equation}
因此未被屏蔽的 't Hooft--Wilson 线构成的群是
\begin{equation}
    \mathsf{D}_1 = \frac{\mathsf{D}_\gamma}{Q_{\mathrm{matter}}} \,,
    \tag{3.204}
\end{equation}
这意味着双荷 1-形式对称性群是
\begin{equation}
    G^{(1)} = \widehat{\mathsf{D}}_1 = \widehat{\dfrac{\mathsf{D}_\gamma}{Q_{\mathrm{matter}}}} \,.
    \tag{3.205}
\end{equation}

\para{离散 theta 项。}$\cG_\gamma$ 理论作用量中额外的离散 theta 角项可以表示为
\begin{equation}
    2\pi \int \cP_\sigma(w_2) \,,
    \tag{3.206}
\end{equation}
其中 $w_2$ 是把 $\cG$ 规范丛提升为 $G$ 规范丛的 $\cZ$ 值 2 次障碍类，而 $\cP_\sigma(w_2)$ 是对 $w_2$ 施行 Pontryagin 平方操作所得的 $\R/\Z$ 值 4 次类。Pontryagin 平方操作由选取一个二次函数 [44] 来指定
\begin{equation}
    \sigma:\ \cZ \to \R/\Z \,,
    \tag{3.207}
\end{equation}
它精化双同态
\begin{equation}
    \eta:\ \cZ \times \cZ \to \R/\Z
    \tag{3.208}
\end{equation}
该双同态对偶于同态 $\gamma$：
\begin{equation}
    \eta(z, z') \equiv (\gamma(z))(z') \in \U(1) \cong \R/\Z \,.
    \tag{3.209}
\end{equation}
这里我们用到了 $\gamma(z) \in \widehat{\cZ}$ 自然地与 $z' \in \cZ$ 缩并给出一个 $\U(1)$ 元素这一事实。

% [Tutor Note: p.58 式 (3.206)--(3.209) 中 "bi-homomorphism" 是双同态：η 在每一变量上分别同态，从而给出 Z 上的 Z 值二次型；σ 是它的二次精化，两者相差一个同态不影响 Pontryagin 平方。这与 (3.194) 的 cochain 公式互为对偶描述。]

\begin{notebox}{二次精化}
从一个双同态 $\eta:\ \cZ \times \cZ \to \R/\Z$ 出发，二次精化 $\sigma:\ \cZ \to \R/\Z$ 是一个二次函数，它满足
\begin{equation}
    \eta(z_1, z_2) = \sigma(z_1 + z_2) - \sigma(z_1) - \sigma(z_2) \,.
    \tag{3.210}
\end{equation}
$\sigma$ 的不同选取在容许自旋结构的 4d 时空上定义了相同的 Pontryagin 平方操作。
\end{notebox}

% [原文 p.59]

\begin{examplebox}{例 3.12：带离散 theta 角的正交辛规范理论}
我们考虑一个与 (3.11) 中所考虑的理论相类似的理论。我们有一个规范代数为
\begin{equation}
    \fg = \fso(4n+2) \oplus \fsp(m)
    \tag{3.211}
\end{equation}
的 4d 规范理论。与之相伴的单连通群是
\begin{equation}
    G = \Spin(4n+2) \times \Sp(m) \,,
    \tag{3.212}
\end{equation}
其中心是
\begin{equation}
    Z(G) = Z(\Spin(4n+2)) \times Z(\Sp(m)) = \Z_4 \times \Z_2 \,.
    \tag{3.213}
\end{equation}
我们取\textbf{规范群}为
\begin{equation}
    \cG = \SO(4n+2) \times \Sp(m) \,,
    \tag{3.214}
\end{equation}
对此我们有
\begin{equation}
    \cZ = \Z_2 \,.
    \tag{3.215}
\end{equation}
我们另外加入处于双基础表示 (3.171) 的物质场，并开启对应于非平凡同态
\begin{equation}
    \gamma:\ \Z_2 \to \Z_2
    \tag{3.216}
\end{equation}
的离散 theta 角。让我们计算 1-形式对称性群。我们有
\begin{equation}
    \cZ \times \widehat{Z}(G) = \Z_2 \times \Z_4 \times \Z_2 \,.
    \tag{3.217}
\end{equation}
未被屏蔽线算符构成的群按 (3.199) 如下构造。把 $\cZ = \Z_2$ 的元素记作 $z \in \{0,1\}$。满同态 $\hat{i}$ 是
\begin{equation}
    \big\{\{(0,0),(2,0),(0,1),(2,1)\} \to \widehat{0},\ \{(1,0),(3,0),(1,1),(3,1)\} \to \widehat{1}\big\} \,.
    \tag{3.218}
\end{equation}
我们应当构造所有配对 $(z, \tilde{\gamma}(z))$。先从 $z = 0 \in \Z_2$ 开始。我们要找 $\widehat{0}$ 的提升，它们是 $\{(0,0),(2,0),(0,1),(2,1)\} \in \widehat{Z}(G) = \Z_4 \times \Z_2$。类似地，从 $z = 1 \in \Z_2$ 出发，我们要找 $\widehat{1}$ 的提升，它们是 $\{(1,0),(3,0),(1,1),(3,1)\} \in \widehat{Z}(G) = \Z_4 \times \Z_2$。于是我们得到
\begin{align}
    \mathsf{D}_\gamma &= \{(0,0,0),(0,2,0),(1,1,0),(1,3,0),(0,0,1),(0,2,1),(1,1,1),(1,3,1)\} , \nonumber \\
    &= \Z_4 \times \Z_2 \subset \cZ \times \widehat{Z}(G) = \Z_2 \times \Z_4 \times \Z_2 \,.
    \tag{3.219}
\end{align}
\end{examplebox}

% [Tutor Note: p.59 式 (3.219) 可逐个复核：z = 0 配 \hat{0} 的 4 个提升，z = 1 配 \hat{1} 的 4 个提升，共 8 个元素。其中 (1,1,1) 的 4 倍为 (0,0,0)、2 倍为 (0,2,0) ≠ 0，故阶为 4；(0,2,1) 的 2 倍为 (0,0,0)，阶为 2；两者生成整个 8 元群，故 D_γ ≅ Z4 × Z2，与 (3.219) 第二行一致。]

% [原文 p.60]

让我们把 $\mathsf{D}_\gamma$ 的 $\Z_4$ 因子的生成元取为 $(1,1,1)$，把 $\mathsf{D}_\gamma$ 的 $\Z_2$ 因子的生成元取为 $(0,2,1)$。物质内容提供的荷是 $(0,2,1)$，这意味着
\begin{equation}
    \mathsf{D}_1 = \Z_4 \,,
    \tag{3.220}
\end{equation}
从而 1-形式对称性群是
\begin{equation}
    G^{(1)} = \Z_4 \,.
    \tag{3.221}
\end{equation}

\section{高阶形式对称性的性质}

在前面的各节中，我们引入了高阶形式对称性，并讨论了各种维数的量子场论中容许这类对称性的许多例子。我们还描述了通过屏蔽对许多一般类别的规范理论中高阶形式对称性的计算。在本节中，我们描述高阶形式对称性的各种性质。特别地，我们将讨论如何把理论与这类对称性的背景场耦合起来。这与 't Hooft 反常有着错综复杂的联系：'t Hooft 反常既是 RG 流的不变量，也刻画了规范化高阶形式对称性的障碍。自 [4] 出现之后，这些性质已被用于各种物理应用。关于这些工作的一个样本，见 [81--121]（排名不分先后）。

\subsection{背景场}

我们从背景场的讨论开始。讨论会略有不同，取决于所讨论的高阶形式对称性是离散的还是连续的。

\para{重温连续的 0-形式对称性。}回忆一下，连续 0-形式对称群 $G^{(0)}$ 的背景场\footnote{人们常常额外加上一个形容词“gauge”，把“background field”称作“background gauge field”。我们这里删去这个形容词，以免与密切相关的规范场概念混淆。在本文中，背景场是时空上的一个固定的、非动力学的场，不在路径积分中求和；而规范场是路径积分中被求和的动力学场。}是时空上 $G^{(0)}$ 丛上的一个联络 $A$，它在局部可以描述为时空上的一个 1-形式场。利用联络的信息，我们可以计算场强 $F$，它是时空上的一个闭 2-形式场。特别地，对于
\begin{equation}
    G^{(0)} = \U(1) \,,
    \tag{4.1}
\end{equation}
我们有
\begin{equation}
    F = \dd A \,.
    \tag{4.2}
\end{equation}

% [Tutor Note: p.60 式 (4.2) 只在局部成立；整体上 F 是闭 2-形式，其 de Rham 类给出第一 Chern 类，而 ∫F/2π ∈ Z 是磁通量子化。这是后文把 U(1) 0-形式对称性的背景场与 2-形式规范场联系起来的起点。]

% [Tutor Note: p.60 脚注 6 的区分是全文反复用到的：背景场 = 固定的、非动力学的、不在路径积分中求和；规范场 = 动力学、要被求和。后文把所有背景场当作探测对称性的外部源，'t Hooft 反常就定义为背景场规范变换下配分函数的相位反常。]
